% !PRINTED 553
\weberpage{553}
%<<CONT
contain a common element must be completely identical, and the division into classes is therefore entirely unambiguous. Every class is completely determined by an arbitrary functional contained in it, a \gesperrt{representative}.

3. \gesperrt{Two functionals associated with one another are also equivalent, and therefore occur in the same class.}

For if $\varphi$ and $\psi$ are associated, then their quotient $\varphi:\psi$ is a unit, and $\varphi$ and $\psi$ are therefore also equivalent.

But a class contains not merely the individual functionals $\varphi$, but all the \gesperrt{ideals} determined by these functionals, and when a more precise designation is necessary we therefore call these classes \gesperrt{functional classes} or, still more frequently, in accordance with the usual usage, \gesperrt{ideal classes}.

4. \gesperrt{The totality of all integral and fractional numbers of the field $\Omega$, combined with all the units and the products of numbers with units, form among themselves a class, which is called the principal class.}

As representative of the principal class one may consider, e.g., the number 1. The principal class, conceived as an ideal class, contains the ideal $\mathfrak{o}$ and will therefore in what follows be denoted by the letter $O$.

5. \gesperrt{In every ideal class there are integral functionals.}

For by the definition, if $\varphi$ is any functional and $\alpha$ a number, $\varphi$ is equivalent to $\alpha\varphi$. But by \S.\,137, 5. we can determine the number $\alpha$, even rational, so that $\alpha\varphi$ becomes an integral functional.

6. \gesperrt{From every ideal class $C$ we can select a representative $\varphi$ which is not only itself an integral functional, but is also relatively prime to an arbitrarily given integral functional $\omega$.}

If, in order to prove this theorem, we first take by 5. an arbitrary integral representative $\varphi$ of the class $C$ and an integer $\alpha$ divisible by $\varphi$, then

\weq{2}{\varphi\chi = \alpha,}
%CONT>>
